% 333_Kfrak_Dualitaet_De_ch.tex
% Dok. 333, FFGFT-Korpus, August 2026 — Kapitel-Datei
% K_frak, Ausrollen und die T~·m-Dualität
% Warum die fraktale Korrektur keine Eigenschaft des Zeitausrollens ist

% ============================================================
\section*{Vorbemerkung}

Im FFGFT-Korpus bedeutet „Ausrollen" bisher stets die
\textbf{Typ-I-Dekompaktifizierung}: Der Messakt wählt den Massenkreis
als Zeit, $S^1_m\to\mathbb{R}_t$, und verwirft dabei die Windungszahl
-- ein verlustbehafteter Schritt (Dok.~285, Dok.~248, Dok.~270).
In Dok.~332 wird daneben ein anderer Vorgang eingeführt: der
\textbf{kanonische Funktionen-Pullback}
$\Phi=p^*:L^2(S^1_m)\to\mathrm{Per}_{R_m}(\mathbb{R}_t)$, der
verlustfrei und exakt ist -- ein reiner Koordinatenwechsel ohne
Informationsverlust.

Dieses Dokument stellt beide Operationen nebeneinander, klärt die Rolle
von $K_{\text{frak}}$ in beiden Fällen, und zeigt, dass die Rekursion
zwingend aus der $\tilde{T}\cdot m$-Dualität folgt.

% ============================================================
\section{Zwei Bedeutungen von „Ausrollen" im Korpus}

\subsection{Typ I: der Messakt (verlustbehaftet)}

Die bisher im Korpus verwendete Bedeutung (Dok.~285 §\,„How FFGFT
unrolls"): Der eingebettete Beobachter wählt den Massenkreis $S^1_m$
als Zeitachse. Die Abbildung
\[
  S^1_m \;\longrightarrow\; \mathbb{R}_t, \qquad
  \tilde T \;\mapsto\; t,
\]
ist \textbf{verlustbehaftet}: Die Windungszahl wird verworfen, das
diskrete Spektrum des kompakten Kreises tritt nicht in $\mathbb{R}_t$
auf. Dies ist keine mathematische Operation der Theorie, sondern die
Brücke zur Beobachtung -- der Torus bleibt die fundamentale Struktur,
$\mathbb{R}_t$ ist das Bild (Dok.~285). Gleichzeitig rollt die
Rekursion $\xi^n$ zu einem kontinuierlichen Dilatationsfluss
$t\mapsto\xi\,t$ mit Log-Periode $|\ln\xi|=\ln 7500\approx 8{,}92$
aus (Dok.~285 §\,„The recursion unrolls along with it").

\subsection{Typ III: der Koordinatenwechsel (verlustfrei)}

Die in Dok.~332 eingeführte Bedeutung: Der kanonische
Funktionen-Pullback
\[
  \Phi = p^* \;:\; L^2(S^1_m) \;\longrightarrow\;
  \mathrm{Per}_{R_m}(\mathbb{R}_t), \qquad
  (\Phi f)(t) = f(t \bmod R_m),
\]
ist eine \textbf{exakte Isometrie} -- kein Informationsverlust, keine
Windungszahl verworfen, Fourier-Basis und Parseval-Identität bleiben
erhalten. Es ist ein Darstellungswechsel innerhalb isomorpher
Hilberträume (R96, Dok.~332). Die verlustbehaftete Richtung liegt
bei $\mathbb{R}_t\to S^1_m$, nicht beim Pullback selbst.

\begin{center}
\begin{tabular}{lp{4.5cm}p{4.5cm}}
\toprule
 & \textbf{Typ I (Messakt)} & \textbf{Typ III (Pullback)} \\
\midrule
Operation & $S^1_m\to\mathbb{R}_t$ & $L^2(S^1_m)\cong\mathrm{Per}_{R_m}(\mathbb{R}_t)$ \\
Verlust & Windungszahl verworfen & keiner \\
Charakter & physikalischer Akt & mathematischer Darstellungswechsel \\
Windungszahl & nicht erhalten & erhalten \\
Informationsverlust & ja & nein \\
Korpus-Belege & Dok.~285, 248, 270 & Dok.~332 (R96) \\
\bottomrule
\end{tabular}
\end{center}

% ============================================================
\section{Die Rolle von $K_{\text{frak}}$ in beiden Fällen}

\subsection{Beim Typ-III-Pullback: unsichtbar}

Ein Laplace-Operator auf dem kompakten $T^4$ ist lokal -- er kann
keine kumulative Größe tragen (Dok.~314). $K_{\text{frak}}$ kommt
daher nicht über den Operator herein, sondern über die
\textbf{Brückenkonstanten} $E_0$, $v$ bei der Skalenumrechnung
(Dok.~314, R72).

Im eingerollten Zustand auf $S^1_m$ ist die Periode
$R_m=\hbar/(mc^2)$ in natürlichen Einheiten exakt -- kein
$K_{\text{frak}}$. Erst wenn $m$ aus der Massenformel eingesetzt wird,
trägt $R_m$ die Korrektur:
\[
  m_i = r_i\cdot\xi_0^{p_i}\cdot v\cdot K_{\text{frak}}
  \quad\Rightarrow\quad
  R_m^{\text{korr.}} = \frac{\hbar}{m_i^{\text{nackt}}\cdot K_{\text{frak}}\cdot c^2}
  = \frac{R_m^{\text{nackt}}}{K_{\text{frak}}}.
\]
$K_{\text{frak}}$ sitzt in der \emph{Wahl von $R_m$}, nicht im
Pullback-Mechanismus. Der Pullback selbst ist in beiden Fällen
(nackt oder korrigiert) exakt.

\subsection{Beim Typ-I-Messakt: sichtbar als Skalenkorrektur}

Beim verlustbehafteten Ausrollen tritt $K_{\text{frak}}$ als
Korrekturfaktor der Brückenkonstanten auf: Das Verhältnis der beiden
$E_0$-Werte (nackt und $\alpha$-verankert) ist $1/\sqrt{K_{\text{frak}}}$
(R72, Dok.~314). Die Korrektur gehört zur \emph{Skalenstruktur} -- zum
Übergang zwischen der kompakten und der beobachtbaren Darstellung --,
nicht zur Zeitgeometrie selbst.

\begin{center}
\begin{tabular}{lp{4.5cm}p{4.5cm}}
\toprule
 & \textbf{Typ I} & \textbf{Typ III} \\
\midrule
$K_{\text{frak}}$ & sichtbar in Brückenkonstanten
  ($E_0$-Verhältnis, Massenformel) & unsichtbar am Isometrie-Schritt;
  sitzt in $R_m$ über $m$ \\
Wo genau & beim SI-Übergang & in der Wahl der Periode \\
\bottomrule
\end{tabular}
\end{center}

% ============================================================
\section{Masse und Zeit rollen gemeinsam}

In der FFGFT-Kompaktifizierung sind $\tilde{T}$ und $m$ eine
\textbf{einzige geometrische Dimension}:
\[
  \tilde{T}\cdot m = 1 \qquad \text{(Grundrelation, Dok.~285)}.
\]
Beide sind keine zwei unabhängigen Größen. Beim Typ-III-Pullback
rollt der Pullback daher $m$ und $t$ \emph{gemeinsam}: Die gesamte
Mode $f\in L^2(S^1_m)$ wird ausgerollt, das Verhältnis $\tilde{T}/m$
bleibt konstant. Die Physik ändert sich nicht, nur die
Koordinatendarstellung.

Ein Ausrollen von $t$ allein -- das $R_m$ fallen lässt -- würde
$\tilde{T}\cdot m=1$ verletzen und ist strukturell unzulässig.
$K_{\text{frak}}$ ist \emph{nicht} der Preis des Ausrollens, sondern
der Preis der \textbf{künstlichen Trennung} von $m$ und $t$ beim
Übergang zu absoluten SI-Einheiten.

% ============================================================
\section{Die Rekursion ist zwingend an die Dualität gebunden}

\subsection{Stabiles Teilchen: Fall B}

Für ein stabiles, isoliertes Teilchen gilt $\xi$ eingefroren
(Fall~B, Dok.~315, Dok.~313): Die Windungszahl ist rational $74/75$,
die Schließung erfolgt nach exakt 75 Umläufen. $K_{\text{frak}}=74/75$
ist die geschlossene Buchhaltung dieser eingerollten Geometrie --
exakt in Bruchrechnung, weil $\xi$ rational gesetzt ist (Dok.~315).

\subsection{Dynamik: Fall C und die Rekursionskette}

Sobald Dynamik ins Spiel kommt -- Wechselwirkung, RG-Lauf,
Massenänderung -- greift eine tiefere Konsequenz von
$\tilde{T}\cdot m=1$: Ändert sich $m$, so ändert sich $\tilde{T}=1/m$
zwingend mit. Die Periode $R_m=\hbar/(mc^2)$ ändert sich, die
Zeitskala der Physik ändert sich, das wirkt auf die effektive Masse
zurück. Die Rekursion (Dok.~295, Dok.~146) lautet:
\[
  \xi_{n+1} = \xi_n\,(1-100\,\xi_n).
\]
Der Schließungsdefekt je Umlauf bleibt nicht konstant, sondern läuft
kumulativ. Die rationale 75-Schließung wird zur logarithmischen Spirale
(Fall~C, Dok.~295):
\[
  \rho = 75\,e^{\beta/2\pi}
  \qquad\text{(äquianguläre Spirale, Selbstähnlichkeit }e\text{)}.
\]

In Fall~C ist $K_{\text{frak}}$ keine Konstante mehr, sondern eine
effektiv skalenabhängige Größe. Die multiplikative Form
$(1-\xi)^{100}$ ist die natürliche Buchhaltung dieses laufenden Falls;
ihr Abstand von $74/75$ -- der Binomialterm $4950\,\xi^2$ (Dok.~315)
-- ist genau der Unterschied zwischen eingefrorener und laufender
Rekursion.

\subsection{Warum Fall A ausgeschlossen ist}

Fall~A ($K_{\text{frak}}=1$, Schließung nach einem Umlauf) würde
bedeuten: $\tilde{T}\cdot m=1$ erzeugt keine Rekursion, weil $m$
sich nie ändert -- ein statisches, wechselwirkungsfreies Universum
mit $D_f=3$ exakt. Das widerspricht der beobachteten Teilchenstruktur
und der Grundanlage des Rahmens.

\subsection{Zusammenfassung der drei Fälle}

\begin{center}
\begin{tabular}{p{1.2cm}p{3.5cm}p{3.5cm}p{3.5cm}}
\toprule
\textbf{Fall} & \textbf{$\xi$} & \textbf{$K_{\text{frak}}$} &
\textbf{Bindung an Dualität} \\
\midrule
A & konstant, $d=0$ & $=1$ & keine Rückkopplung, ausgeschlossen \\
B & eingefroren & $=74/75$ exakt & stabile Mode, Dualität ruhend \\
C & läuft (Rekursion) & skalenabhängig & Dualität erzeugt laufende
  Rückkopplung $m\leftrightarrow\tilde{T}$ \\
\bottomrule
\end{tabular}
\end{center}

$K_{\text{frak}}=74/75$ ist die \textbf{Projektion von Fall~C auf
Fall~B}: exakt in der eingerollten, statischen Beschreibung; erste
Näherung in der dynamischen.

% ============================================================
\section*{Zusammenfassung}

\begin{quote}
\emph{„Ausrollen" bedeutet im FFGFT-Korpus bisher stets den
verlustbehafteten Messakt (Typ~I, Dok.~285): Windungszahl verworfen,
$S^1_m\to\mathbb{R}_t$. Der Funktionen-Pullback aus Dok.~332 ist eine
andere, verlustfreie Operation (Typ~III): ein exakter
Darstellungswechsel, der $m$ und $t$ gemeinsam ausrollt und an dem
$K_{\text{frak}}$ nicht beteiligt ist. $K_{\text{frak}}$ sitzt in der
Wahl der Periode $R_m$ über die Massenformel und tritt sichtbar erst
beim Übergang zu absoluten SI-Einheiten auf. Die Rekursion
$\xi_{n+1}=\xi_n(1-100\,\xi_n)$ ist zwingend an $\tilde{T}\cdot m=1$
gebunden: Jede Massenänderung zieht eine Zeitänderung nach sich, was
rekursiv auf die Masse zurückwirkt. Fall~B (stabiles Teilchen,
$K_{\text{frak}}=74/75$) ist die Näherung für stillgestellte Dualität;
Fall~C (laufende Rekursion, Log-Spirale) ist die vollständige
dynamische Beschreibung."}
\end{quote}

\medskip\noindent
\textbf{Verweise:} Dok.~133 ($K_{\text{frak}}$-Herleitung),
Dok.~146 (Rekursion), Dok.~248 (Beobachter),
Dok.~270 (Typ-I-Theorem), Dok.~285 (Dekompaktifizierung),
Dok.~295 (Log-Spirale), Dok.~313 (drei Fälle),
Dok.~314 (lokaler Operator, Brückenkonstanten),
Dok.~315 (Formfrage $K_{\text{frak}}$),
Dok.~332 (Pullback-Isometrie, R96).
